% Canonical Paper IV section.
The projective operator in Section~\ref{sec:history-equality} is already a
condensation of a marked history.  We now introduce a second, algebraic
condensation: an operator is replaced by the image--kernel pair of a rank-one
projector.  No topology, norm, probability, or limiting dynamics is used in
this section.

\subsection{Points, copoints, and regular bivaluations}

Let $V=K^2$ and let $V^\vee$ be its algebraic dual.  A nonzero vector $v$
defines a point $[v]\in\PP(V)$, while a nonzero covector $\varphi$ defines the
projective line $\PP(\ker\varphi)\in\PP(V)$.  The pair is transverse when
$\varphi(v)\ne0$.

\begin{definition}[Regular bivaluation]
A regular bivaluation over $K$ is an ordered pair
\[
  (q^+,q^-)\in\PP^1(K)\times\PP^1(K),
  \qquad q^+\ne q^-.
\]
The first member is called the image direction and the second the kernel
direction.  We write
\[
  \Biv_K^{\mathrm{reg}}
  \defeq \PP^1(K)\times\PP^1(K)\setminus\Delta.
\]
\end{definition}

The order matters.  Interchanging image and kernel replaces a projector by its
complement.

\begin{theorem}[Bivaluation--projector correspondence]
\label{thm:bivaluation-projector}
There is a natural bijection
\[
  \Biv_K^{\mathrm{reg}}
  \longleftrightarrow
  \{\Pi\in\End(V):\Pi^2=\Pi,\ \rank\Pi=1\}.
\]
If $q^+=[v]$, $q^-=\PP(\ker\varphi)$, and $\varphi(v)\ne0$, the
corresponding projector is
\begin{equation}
  \Pi_{v,\varphi}=\frac{v\otimes\varphi}{\varphi(v)}.
  \label{eq:rank-one-projector}
\end{equation}
\end{theorem}

\begin{proof}
For transverse lines $L^+$ and $L^-$ one has $V=L^+\oplus L^-$.  There is a
unique linear projection onto $L^+$ along $L^-$.  Formula
\eqref{eq:rank-one-projector} gives it and is unchanged if $v$ or $\varphi$
is rescaled.  Conversely, an idempotent splits
$V=\im\Pi\oplus\ker\Pi$.  Rank one makes both summands distinct projective
lines.  The two constructions are inverse.
\end{proof}

In the affine chart
\[
  v_a=\begin{pmatrix}-a\\1\end{pmatrix},
  \qquad
  \varphi_b=\begin{pmatrix}-b&1\end{pmatrix},
\]
one has $\varphi_b(v_a)=1+ab$ and hence
\begin{equation}
  \Pi_{a,b}
  =\frac{1}{1+ab}
    \begin{pmatrix}ab&-a\\-b&1\end{pmatrix}.
  \label{eq:local-projector}
\end{equation}
The divisor $1+ab=0$ is a genuine incidence failure of the transverse
splitting.  It must not be confused with the pole of one particular coordinate
chart.

\subsection{The homogeneous-space tower}

Fix a reference ordered pair $(q_0^+,q_0^-)$ and set
\[
  G=\PGL_2(K),
  \qquad
  H=\Stab_G(q_0^+,q_0^-),
\]
\[
  B_+=\Stab_G(q_0^+),
  \qquad
  B_-=\Stab_G(q_0^-).
\]
Then $H=B_+\cap B_-$.  We use right cosets in $G/H$ and $G/B_\pm$.

\begin{theorem}[Projective quotient tower]
\label{thm:projective-quotient-tower}
There are natural identifications
\[
  G/H\cong\Biv_K^{\mathrm{reg}},
  \qquad
  G/B_+\cong\PP^1(K),
  \qquad
  G/B_-\cong\PP^1(K).
\]
Under these identifications, the maps induced by
$H\subseteq B_\pm$ forget respectively the kernel or image member of the
ordered pair.
\end{theorem}

\begin{proof}
The group $G$ acts transitively on ordered pairs of distinct projective points.
The stabilizer of the reference pair is exactly $H$, so orbit--stabilizer gives
$G/H$.  The same argument for a single projective point gives the two Borel
quotients.  The descriptions of the forgetful maps follow from the subgroup
inclusions.
\end{proof}

Combining this result with projective evaluation gives the central algebraic
tower
\begin{equation}
\boxed{
  \Hist_K
  \xrightarrow{\ \rho\ }
  G
  \xrightarrow{\ \pi_H\ }
  G/H\cong\Biv_K^{\mathrm{reg}}
  \longrightarrow
  G/B_\pm\cong\PP^1(K).
}
\label{eq:paper4-main-tower}
\end{equation}
The first fiber consists of literal histories inducing one operator.  The
second consists of projective frames over one image--kernel result.  These are
different kinds of forgotten data.

\subsection{A third point restores the projective frame}

\begin{definition}[Framed bivaluation]
A framed regular bivaluation is a pairwise distinct ordered triple
$(q^+,q^-,q^\infty)$ of points of $\PP^1(K)$.  The first two determine the
rank-one projector; the third selects a projective frame above it.
\end{definition}

\begin{theorem}[Frame torsor]
\label{thm:frame-torsor}
After a reference triple has been chosen, the set of framed regular
bivaluations is a $G$-torsor.  The projection
\[
  (q^+,q^-,q^\infty)\longmapsto(q^+,q^-)
\]
has fibers on which $H$ acts freely and transitively.  Equivalently,
$G\to G/H$ is an $H$-torsor.
\end{theorem}

\begin{proof}
A projective transformation is uniquely determined by its values on three
distinct projective points.  Hence $G$ acts simply transitively on ordered
projective frames.  Holding the first two points fixed leaves exactly their
stabilizer $H$, which acts simply transitively on the allowed third points.
\end{proof}

For a Lie group this torsor may be equipped with a connection, but no such
connection is canonical over an arbitrary field.  A coordinate in $H$ also
requires a reference lift or local section.

\begin{proposition}[Stabilizer-valued operator residue]
\label{prop:stabilizer-residue}
Let $\Pi_0$ be the projector of the reference bivaluation and define
$\pi_H(g)=g\Pi_0g^{-1}$.  Then
\[
  \pi_H(g_1)=\pi_H(g_2)
  \quad\Longleftrightarrow\quad
  g_2^{-1}g_1\in H.
\]
Thus two operator-level processes with the same condensed projector differ by
an $H$-valued frame residue.
\end{proposition}

\begin{proof}
Equality of the two conjugated projectors is equivalent to
\[
  (g_2^{-1}g_1)\Pi_0(g_2^{-1}g_1)^{-1}=\Pi_0.
\]
The stabilizer of a rank-one projector is precisely the stabilizer of its
ordered image and kernel lines, namely $H$.
\end{proof}

A static idempotent already splits $V=\im\Pi\oplus\ker\Pi$.  It therefore does
not carry a nontrivial linear extension class by itself.  Process information
discarded by $G\to G/H$ belongs instead to a lift, a connection and its
holonomy, or some additional decoration of the history.

\subsection{Finite-field calibration}

\begin{proposition}[Exact finite-field fiber counts]
\label{prop:finite-field-counts}
For $K=\F_q$,
\[
  |\PP^1(\F_q)|=q+1,
  \qquad
  |\Biv_{\F_q}^{\mathrm{reg}}|=q(q+1),
\]
\[
  |\PGL_2(\F_q)|=q(q^2-1),
  \qquad
  |H|=q-1.
\]
Consequently every fiber of
$\PGL_2(\F_q)\to\Biv_{\F_q}^{\mathrm{reg}}$ has $q-1$ elements.
\end{proposition}

\begin{proof}
There are $q+1$ projective points and then $q$ choices for a distinct second
point.  Moreover
$|\GL_2(\F_q)|=(q^2-1)(q^2-q)$; quotienting by the $q-1$ nonzero scalar
matrices gives $q(q^2-1)$.  The stabilizer of $(0,\infty)$ consists of
$z\mapsto az$ with $a\in\F_q^\times$, and hence has order $q-1$.
\end{proof}

\begin{example}[The field $\F_7$]
One obtains
\[
  |G|=336,
  \qquad |G/H|=56,
  \qquad |H|=6.
\]
The 336 ordered projective frames condense to 56 ordered image--kernel pairs,
with six frame lifts above each pair.
\end{example}

\begin{corollary}[Uniform finite-frame information]
\label{cor:finite-frame-entropy}
Let $G=\PGL_2(\F_q)$ be uniformly distributed and let $Q=gH$.  Then
\[
  H(g\mid Q)=\log_2(q-1).
\]
This is the exact conditional information in the uniformly distributed frame
coordinate above a bivaluation.
\end{corollary}

\begin{proof}
Every quotient fiber has $q-1$ elements and the conditional distribution on
that fiber is uniform.
\end{proof}

The corollary is an algebraic information statement.  Whether these bits are
needed by an online computation depends on the continuation experiment of the
next section.

\subsection{Algebraic and dynamical condensation}

\begin{definition}[Algebraic condensation]
Algebraic condensation is passage to an idempotent or quotient state, such as
$g\mapsto gH$ or $g\mapsto g\Pi_0g^{-1}$.  It requires no limiting sequence.
\end{definition}

\begin{definition}[Dynamical condensation]
Dynamical condensation is convergence, after a declared normalization, of a
sequence of linear lifts to a singular rank-one operator, for example
\[
  \frac{\widetilde g_n\cdots\widetilde g_1}{r_n}
  \longrightarrow \Pi
  \quad\text{in }\End(V).
\]
It requires a topology or valuation, a choice of lifts, and a normalization.
\end{definition}

\begin{proposition}[No automatic dynamical condensation over a finite field]
In the discrete topology on $\PGL_2(\F_q)$, every sequence of powers is
periodic and convergence occurs only when the sequence is eventually constant.
Thus the existence of rank-one idempotents over $\F_q$ does not imply a
nontrivial attracting limit in $G$.
\end{proposition}

The distinction will remain active throughout the paper.  The quotient
$G/H$ is an exact algebraic state space; whether it is sufficient for future
computation is a separate question.
